\section*{Document status and decision record}
\addcontentsline{toc}{section}{Document status and decision record}
\statusnote{Design specification; not an implementation report}{This document
defines requirements for a future station network. It does not claim that the
released v0a application implements these services or has passed the tests below.}

The agreed deployment comprises ten computers on one reachable local network:
Launch Telemetry, Launch Video, and Telemetry/Video computers at each of four
Away stations. The user has accepted the embedded-service architecture and H.264
video. This revision makes the proposed network protocol precise enough to guide
implementation and interoperability tests. Transport defaults, rates and resource
limits below are engineering proposals, to be qualified on the actual network.

\begin{longtable}{L{35mm}L{126mm}}
\toprule Decision & Baseline in this revision\\\midrule
\endfirsthead
\toprule Decision & Baseline in this revision\\\midrule
\endhead
Discovery & Bonjour-compatible mDNS and DNS-SD; one advertised application service per computer. Manual endpoint fallback.\\
Metadata and snapshots & HTTP API with JSON; current values come from memory, not a repeatedly rewritten text file.\\
Live telemetry and health & WebSocket event stream, independent of video, with source identity, sequence numbers, bounded resumption and explicit gaps.\\
Video & H.264/AVC, MPEG transport stream container, SRT live transport over UDP. Preview and main renditions.\\
Station naming & UI says Launch station. Requested network aliases retain \code{rt-base-*}; Away aliases use \code{rt-gs1-*} through \code{rt-gs4-*}.\\
Hardware commands & Existing serial protocols remain separate. Unspecified uplink/target-switch/remote-control commands are unsupported in protocol 1.0.\\
Reference application & The released v0a/0.1.1 manual remains historical evidence of implemented behavior. This specification is a forward design.\\
\bottomrule
\end{longtable}

\textbf{Reading convention.} SHALL is required to claim conformance to this
proposed protocol; SHOULD permits a documented engineering exception; MAY is
optional. These words define requirements, not a claim of user approval of every
numeric default. All times are milliseconds unless a unit is stated. JSON field
names and enum values are case-sensitive. Source listings are illustrative where
identified; complete machine-readable examples accompany this report.

\textbf{Scope.} Discovery, identity, mission association, monitoring, raw telemetry
provenance, video transport, reconnection, UI behavior, limits and verification.
The design runs locally without a cloud service or Internet access. It does not
specify a navigation estimator, board firmware, LTU management API, RF modulation,
voice communications or launch authorization procedure.

\textbf{Revision history.} Revision 0.1, 25 October 2026: first requirements and
protocol draft following the ten-computer network discussion. No application
source changes are part of this document.
